\scanpage{72}
\uline{Application} Toutes les constructions (``universelles'') dans $P_*^{\unsure{2|\Sigma}}$ \deleted{qui} (ou $\deleted{\unsure{P}}Q_*^{\unsure{1}},P_{0*},\mathbb{K}_*$) qui se font \unsure{uniquement|intrinsèquement} en termes de sa structure $k$-lin et sa $\otimes$-structure, \deleted{donnent} \inserted{se reflètent en} des constructions analogues sur un $C^{\bullet}\in\mathrm{ob}\,B^{\bullet}$ muni de structure multiplicative.

\uline{Ex 1} \deleted{\unsure{1 $\mathbb{C}^{\bullet}_*$ On a un hom|$\Psi^0$ $\mathbb{C}^{\bullet}_*$ On a un hom}} \inserted{Comme} $\mathbb{C}^{\bullet}_*\otimes \mathbb{C}^{\bullet}_*$ est une résolution $\Sigma^{\unsure{*}}_*$-inj. de $\Psi^0_{\unsure{\bullet|0}*}$ dans $P_{0*}$ tt comme $\mathbb{C}^{\bullet}_*$ lui-même on trouve
\[
\underset{\substack{\downarrow\\\unsure{\substack{\text{complexe}\\\text{simple associé}}|\substack{\text{avec les}\\\text{signes usuels}}}}}{\operatorname{Tot}^{\bullet}}(\mathbb{C}^{\bullet}_*\deleted{\unsure{\ast}}\otimes \mathbb{C}^{\bullet}_*)\longrightarrow \mathbb{C}^{\bullet}_*\quad\text{défini à homotopie}
\]
près, d'où un hom
\diagram{\begin{tikzcd}[ampersand replacement=\&,column sep=large,row sep=small]
F_{C^{\bullet}}(\mathbb{C}^{\bullet}_*\deleted{\unsure{\otimes}}*\mathbb{C}^{\bullet}_*)\arrow[r] \& F(\mathbb{C}^{\bullet}_*)\arrow[d,equal]\quad\text{défini à hom. près}\\
F_{C^{\bullet}}(\mathbb{C}^{\bullet})\otimes_B F_{C^{\bullet}}(\mathbb{C}^{\bullet})\arrow[u] \& C^{\bullet}\\
\mathbb{C}^{\bullet}\otimes C^{\bullet}\arrow[u,equal] \&
\end{tikzcd}}
d'où en composant un hom défini à homotopie près
\[
C^{\bullet}\otimes_B C^{\bullet}\longrightarrow C^{\bullet}
\]
avec conditions d'associativité, de comm. (attention aux signes!) et d'unité -- le tout à homotopie près. Il en résulte une structure de ``$\otimes B$ Alg. \unsure{graduée|gradué}'' sur $H^*(C^{\bullet})$, avec les propriétés habituelles.
% archive sheet number: 71
